\section{Necessity Derivation of G1--G6 from A1}
\label{sec:g1-g6-necessity-derivation}

The trace-to-scheme export theorem of \S\ref{sec:trace-to-scheme-export-v35} states that an APF trace value may be compared to a physical observable only when the six gates G1--G6 are satisfied. The empirical training set (the six mass routes of 2026-05-10 plus the three closed Reconstruction Programs walked through the gates) confirms that the six-gate pattern discriminates route-by-route status correctly. This section establishes a stronger claim: \emph{the six gates are necessary conditions, derivable from APF's foundational axiom $A_1$.} Each gate has the structural form of a single $A_1$ commitment applied to one moment of the export operation.

Sufficiency of $G_1$--$G_6$ (no missing $G_7$) and uniqueness of the route-status determination are separate theorem programs not addressed here.

\subsection{Foundational commitments and categorical setup}
\label{subsec:setup}

\paragraph{Foundational commitments.} $A_1$: every causally connected region of admissibility space has finite capacity, and every physical distinction carries positive enforcement cost. MD (corpus-derived from $A_1$ + BW): every physical distinction carries a uniform floor $\varepsilon^{*} > 0$. BW (corpus-derived from $A_1$ + finite-physical-regime): the cost spectrum is non-degenerate and bounded. $A_2$: the finite-physical-regime hypothesis $C_\Gamma < \infty$.

\paragraph{Categorical structure.} The \emph{trace category} $\mathcal{T}$ is the interface where APF natively derives, with substrate $S_T$, distinction set $\mathfrak{D}_T$, and capacity $C_T$. A physical-scheme category $\mathcal{S}_i$ (OS, $\overline{\rm MS}$, MSR, pole, MC, lattice $\overline{\rm MS}$(2 GeV), \dots) is an external interface with its own enforcement triple $(S_{\mathcal{S}_i}, \mathfrak{D}_{\mathcal{S}_i}, C_{\mathcal{S}_i})$. A \emph{transport map} $\tau: \mathfrak{D}_T \to \mathfrak{D}_{\mathcal{S}_i}$ takes a trace value $T \in \mathfrak{D}_T$ to an element $\tau(T) \in \mathfrak{D}_{\mathcal{S}_i}$ that can be compared to a measurement $M_{\mathcal{S}_i}$ in the same category.

The six gates $G_1$--$G_6$ are conditions on the pair $(\mathcal{S}_i, \tau)$. We prove each is necessary in turn.

\subsection{Lemma G1: codomain identification}
\label{subsec:lemma-G1}

\begin{lemma}[Codomain Identification]
\label{lem:G1}
Any admissible export must declare its target codomain category $\mathcal{S}$.
\end{lemma}

\begin{proof}
The trace category $\mathcal{T}$ and any physical-scheme category $\mathcal{S}_i$ are distinct interfaces. By the FD1 enforcement-substrate triple structure, each interface owns its own distinction set, and elements of distinct interfaces are not directly comparable --- comparison requires identifying the interface in which it takes place.

Suppose the export attempts to compare $T \in \mathfrak{D}_T$ to ``the measurement $M$'' without specifying which $\mathcal{S}_i$ the measurement lives in. Then $M \in \mathfrak{D}_?$ for an unspecified interface. The measurement's distinction-cost (positive in \emph{its} interface by $A_1$) is not located in any specific cost-budget. The comparison $T \stackrel{?}{=} M$ proceeds with one side ($T$) having a paid cost in $\mathfrak{D}_T$ and the other side ($M$) having its cost suspended.

$A_1$ forbids the suspension: every physical distinction has positive cost paid in \emph{some} interface. Comparing across an unidentified interface conflates distinctions whose costs are paid in different ledgers. Therefore $\mathcal{S}_i$ must be declared. $\square$
\end{proof}

The structural content $G_1$ enforces: scheme conversion is a real distinction (numerically: $M_W^{\rm OS} \ne M_W^{\overline{\rm MS}} \ne M_W^{\rm pole}$), and that distinction must be paid for by being declared.

\subsection{Lemma G2: transport map existence and non-inversion}
\label{subsec:lemma-G2}

\begin{lemma}[Transport Map Existence]
\label{lem:G2}
Any admissible export must supply a transport map $\tau: \mathfrak{D}_T \to \mathfrak{D}_{\mathcal{S}}$ that is either non-inverse or an identity admission of $\mathcal{S} = \mathcal{T}$.
\end{lemma}

\begin{proof}
Given $G_1$ (codomain declared), the comparison requires an element of $\mathfrak{D}_{\mathcal{S}}$. The trace value $T \in \mathfrak{D}_T$. If $\mathfrak{D}_T \ne \mathfrak{D}_{\mathcal{S}}$ as distinction sets --- i.e., $\mathcal{S}$ is genuinely external to the trace category --- then a map $\tau$ producing an element of $\mathfrak{D}_{\mathcal{S}}$ from $T$ is required for the comparison to type-check.

Two ways the gate can close:

\paragraph{(a) Non-inverse map.} $\tau$ is a function constructed from the framework's structural content (transport equations, scale conversions, RG running, scheme matching) that takes $T$ and produces $\tau(T) \in \mathfrak{D}_{\mathcal{S}}$ without using $M_{\mathcal{S}}$ as input. The cost of producing $\tau$ is the cost of the structural distinctions it encodes.

\paragraph{(b) Identity codomain admission.} $\mathcal{S} = \mathcal{T}$ structurally: the codomain is a category APF natively populates. The bottom self-scale case ($m_b(m_b)_{\overline{\rm MS}}$, scheme scale set by the trace value itself) admits this reading. No transport is required because no inter-category jump occurs.

Suppose neither: $\tau$ is supplied as the inverse of some forward map $\mu: \mathfrak{D}_{\mathcal{S}} \to \mathfrak{D}_T$ that consumed $M_{\mathcal{S}}$. Then $\tau(T) = \mu^{-1}(T)$ produces, by construction of $\mu$ from $M_{\mathcal{S}}$, the value $M_{\mathcal{S}}$ itself. The export then claims $\tau(T) = M_{\mathcal{S}}$ --- but this is not a derivation; it's the function $\mu$ being used as its own inverse, with the result given to it. By $A_1$ (developed in $G_5$), zero distinction-cost was paid; therefore no distinction was derived; therefore no admissible export.

Hence $G_2$ must close as (a) or (b). $\square$
\end{proof}

\subsection{Lemma G3: external-input disclosure}
\label{subsec:lemma-G3}

\begin{lemma}[Constants Ledger]
\label{lem:G3}
Every external constant or scale used in $\tau$ must be declared in a constants ledger.
\end{lemma}

\begin{proof}
Let $\tau(T; c_1, \dots, c_n) \in \mathfrak{D}_{\mathcal{S}}$ be the transport map parameterized by constants $\{c_i\}$. By $A_1$, every distinction APF claims to derive carries cost paid by APF's enforcement budget at the trace interface. External constants ($\alpha$, $G_F$, $M_Z$, $\alpha_s(M_Z)$, lattice values, PDG averages) live in \emph{other} interfaces; their cost was paid by other ledgers.

Suppose $c_j$ is used in $\tau$ but not declared. The export claim presents $\tau(T)$ as APF content --- derivable from $T$ via the framework's structure. But $\tau$ uses $c_j$, whose cost APF did not pay. The export conflates two cost-budgets: APF's (covering $T$ and the structural form of $\tau$) and the external interface's (covering $c_j$).

$A_1$'s downstream commitment: derivable content carries the framework's paid cost; external content does not. Conflating them is overclaim --- APF presents content as derived when part of it was input. Therefore $\{c_i\}$ must be declared. $\square$
\end{proof}

\subsection{Lemma G4: uncertainty propagation}
\label{subsec:lemma-G4}

\begin{lemma}[Uncertainty Propagation]
\label{lem:G4}
The uncertainty on the trace value $T$ and the constants ledger $\{c_i\}$ must be propagated through $\tau$ to produce a covariance on $\tau(T)$.
\end{lemma}

\begin{proof}
By BW, the cost spectrum at every interface is non-degenerate and bounded; equivalently, distinctions live with finite spread, not at the cost floor. The trace value $T$ has spread $\sigma_T \ge \varepsilon^{*} > 0$ (from MD); each constant $c_i$ has spread $\sigma_{c_i}$ from its source-interface's measurement or derivation uncertainty.

Suppose $\tau(T)$ is reported with zero spread, or with spread that ignores the input spreads. Then the export effectively claims $\tau$ collapsed all input distinctions to a single point in $\mathfrak{D}_{\mathcal{S}}$. By BW, transport maps generically do not collapse spreads --- they propagate them.

The contrapositive: if zero spread is reported on $\tau(T)$ when the inputs have positive spread, then either (i) $\tau$ is constant over the input distribution (no information transferred --- vacuous transport) or (ii) the input spreads were ignored (cost was hidden). Both options violate $A_1$. Therefore the spread propagation must be reported. $\square$
\end{proof}

For a structural-codomain route ($\mathcal{S}$ a theorem-shape rather than a number-shape), Lemma~\ref{lem:G4} reads in its structural form: the \emph{domain of validity} of the structural closure is declared, with parked or excluded regimes explicit (cf.\ \S\ref{sec:structural-codomain-extension}). The numerical-form covariance and the structural-form domain of validity are both expressions of the same $A_1$ commitment: every distinction's cost must be accounted for.

\subsection{Lemma G5: no-target-as-input}
\label{subsec:lemma-G5}

\begin{lemma}[No-Smuggling]
\label{lem:G5}
The target observable $M_{\mathcal{S}}$ must not appear as an input to $\tau$ or to the trace-side derivation of $T$.
\end{lemma}

\begin{proof}
This is the contrapositive of $A_1$ applied to derivation. $A_1$ says distinctions cost positive enforcement. A derivation is a chain of distinction-enforcing steps with positive total cost.

Suppose $M_{\mathcal{S}}$ appears as input. Decompose $\tau(T) = M_{\mathcal{S}} + \rho$ where $\rho$ is whatever non-target content the route produced. The cost paid by APF for the export is the cost of producing $\rho$, not the cost of $\tau(T) = M_{\mathcal{S}}$. The distinction ``$\tau(T) = M_{\mathcal{S}}$'' was \emph{given}, not derived; APF paid no enforcement cost for that distinction.

But the export claim is ``we derived $\tau(T) = M_{\mathcal{S}}$.'' By $A_1$, derivation requires cost; if the cost paid is only $\rho$'s, the derivation is only of $\rho$. Claiming derivation of $\tau(T) = M_{\mathcal{S}}$ when only $\rho$ was paid for is overclaim.

Equivalently: no enforcement cost paid means no distinction derived. If $M_{\mathcal{S}}$ is consumed as input, the distinction-content of the export collapses to $\rho$ alone; the rest is a tautology. $A_1$ forbids presenting tautology as derivation. $\square$
\end{proof}

\paragraph{Audit-first refusals as $G_5$ instances.} The framework's audit-first quarantine of fits maps directly to $G_5$ violations. The v18.0 W-route numeric projection over the v17.0 tensor scaffold consumed the counterterm target as a fit objective; under Lemma~\ref{lem:G5} the fit fails $G_5$ and is correctly bank-registered as quarantined-non-physical. The v27 charged-lepton three-mode reconstruction would consume $(m_e, m_\mu, m_\tau)$ as fit parameters; quarantined as target-fitted under $G_5$. The v32 light-quark uniform scaling would consume PDG $\overline{\rm MS}$(2 GeV) values as scale-calibration targets; refused under $G_5$. Each refusal is now theorem-required, not a stylistic preference.

\subsection{Lemma G6: residual assignment}
\label{subsec:lemma-G6}

\begin{lemma}[Residual Channel]
\label{lem:G6}
The residual $\rho = \tau(T) - M_{\mathcal{S}}$ must be assigned to a named channel.
\end{lemma}

\begin{proof}
After the comparison, $\rho \in \mathfrak{D}_{\mathcal{S}}$ is a real distinction --- it distinguishes $\tau(T)$ from $M_{\mathcal{S}}$ in the same category. By $A_1$, every element of $\mathfrak{D}_{\mathcal{S}}$ has positive cost paid in \emph{some} interface.

Suppose $\rho$ is left uninterpreted. The framework has produced content (the comparison result) including a residual whose cost is paid by no interface. This is exactly the configuration $A_1$ forbids: a distinction in the framework's content with no paying ledger.

The four admissible channels (each pays the cost in a different ledger):
\begin{enumerate}
\item \textbf{Propagated covariance.} If $|\rho| \lesssim \sigma_{\tau(T)}$ from $G_4$, the residual is consistent with the input/theory-uncertainty distribution. Cost paid by the covariance budget; export reads as $[P_{\rm export\ candidate}]$.
\item \textbf{Named operator-of-this-route the framework owes.} If $\rho$ exceeds the propagated covariance, it can be assigned to a named structural object whose derivation is open --- e.g., the charged-lepton generation residual operator (v27 case). Cost is to-be-paid by a future framework theorem.
\item \textbf{Named scheme-convention difference.} If $\rho$ is the difference between two declared scheme conventions both admitted by external interfaces, the cost is paid by the scheme-conversion bookkeeping at the external interface.
\item \textbf{Named theory-uncertainty band of an admitted external evaluator.} If $\tau$ relies on an external evaluator (DIZET in the W route), $\rho$ within that evaluator's published theory-uncertainty band is paid by the evaluator's external ledger.
\end{enumerate}

If none of (1)--(4) hold, $\rho$ is uninterpreted, $A_1$ is violated, and the export is inadmissible. Therefore $\rho$ must be channel-assigned. $\square$
\end{proof}

For a structural-codomain route, Lemma~\ref{lem:G6} reads in its structural form: \emph{edge cases at the derivation's boundary are named}. Each named edge case has an explicit parking destination (cf.\ \S\ref{sec:structural-codomain-extension}, RP1 walk: quantum-gravity regime parked as RP5 long-term capstone; far-from-equilibrium named in Paper 25 H4).

\subsection{Closing observations}
\label{subsec:closing-g1g6}

\paragraph{What this derivation establishes.} $G_1$--$G_6$ are necessary conditions on any trace-to-scheme transport, derivable from $A_1$ + MD + BW + the trace/scheme category distinction. They are not arbitrary engineering choices; each gate is the structural form of a single $A_1$ commitment applied to one moment of the export operation. The empirical six-for-six (mass routes) plus three-for-three (closed RPs) discrimination across 2026-05-10's articulated routes is now backed by foundational derivation: the routes land where the gates put them because the gates are $A_1$'s structural form on the export operation.

\paragraph{What this derivation does not establish.} Sufficiency of $G_1$--$G_6$ --- that closing all six gates \emph{guarantees} admissible export --- remains open. A sufficiency proof would proceed by structural exhaustion of the trace-to-scheme operation's surfaces: any additional gate $G_7$ would correspond to a structural surface not in the list (codomain identification, transport map, input declaration, uncertainty propagation, target-as-input, residual handling), and no such surface exists. The exhaustion proof is the natural next theorem program after this artifact lands. Likewise, uniqueness of the route-status determination (each gate-availability pattern leads to exactly one named status) remains conjectural pending deductive proof.

\paragraph{Practical content.} Three things become structural rather than methodological once the necessity proof lands. First, the framework's audit-first refusals of fits (v18.0 W-route, v27 charged-lepton, v32 light-quark) are theorem-required, not stylistic preferences. Under $A_1$, the framework had no choice but to refuse. Second, the route-by-route epistemic specificity that emerged on 2026-05-10 is now derived from which subset of $G_1$--$G_6$ each route closes, not just observed empirically. Third, the next theorem program is named precisely: derive sufficiency by structural exhaustion; derive uniqueness of status determination from the gate-pattern combinatorics.
